\documentstyle[% 


righttag% 


,11pt% 


,amssymb% 


,russian%               remove in Eng. 


,izvru%                 Set izven% in Eng. 


]{amsart} 


 


%       Math definitions 


 


\newcommand{\tr}{\operatorname{tr}} 


 


\theoremstyle{plain} 


\theorembodyfont{\it} 


\newtheorem{lemnonum}{\hskip\parindent Лемма} 


\newtheorem{thmnonum}{\hskip\parindent Теорема}


\renewcommand{\thelemnonum}{} 


\renewcommand{\thethmnonum}{}


 


%\hoffset=-15mm%                  For Eng. 


 


\setcounter{page}{29} 


\god{1998} 


\nomer{2~(429)} %                Remove in Eng. 


%\nomer{Vol.\,42, No.\,2}%        Set in Eng. 


% \str{75--81}%                    Set in Eng. 


\udk{519.216} 


\author{\it Ю.В.\,Зайцева} 


% NB:   ^^ remove in Eng.


\title{Управляемые линейные стохастические системы\\


с коэффициентами, зависящими от случайной последовательности} 


 


%\thanks{Работа выполнена при поддержке РФФИ, грант \No. 94-01-00183.} 


\date{} 


% \pagestyle{myheadings}%         Set in Eng. 


 


\pagestyle{plain} %              Remove in Eng. 


\begin{document} 


 


% \thispagestyle{plain}%          Set in Eng. 


 


\maketitle 


 





В работе исследуется управляемое линейное стохастическое


рекуррентное уравнение с коэффициентами, зависящими от


последовательности случайных элементов с известным совместным


распределением.


Рассматривается задача оптимального управления


относительно квадратичного функционала


стоимости.





Подобные системы


изучались в [1] и [2]. В [1] исследуется линейное стохастическое


рекуррентное уравнение с коэффициентами, зависящими от марковской


цепи.





Используемый в данной статье метод построения оптимального управления


отличается от


методов указанных работ. При этом оптимальное управление системой,


содержащей марковскую цепь, получено как частный


случай оптимального управления системой, зависящей от произвольной


последовательности случайных элементов.


Кроме того, в отличие от [1],


стохастические члены рассматриваемого уравнения могут зависеть от


решения и управления.





Пусть $(\Omega,F,P)$ --- полное вероятностное пространство;


$\bigl\{\theta_{k},\; k=\overline{0,N-1}\bigr\}$ --- последовательность


случайных элементов со значениями в измеримом пространcтве $(Z,\Psi)$;


$\bigl\{w_{ik},\;k=\overline{0,N-1}\bigr\}$ --- последовательность


независимых


$p_{i}$-мерных случайных векторов с $Mw_{ik}=0$ и


$Mw_{ik}w_{ik}^{*}=I_{p_{i}}$,


$i=1,2,3$;


случайные последовательности $w_{1}=\bigl\{w_{1k},\;


k=\overline{0,N-1}\bigr\}$,


$w_{2}=\bigl\{w_{2k},\; k=\overline{0,N-1}\bigr\}$,


$w_{3}=\bigl\{w_{3k},\; k=\overline{0,N-1}\bigr\}$ и


$\theta=\bigl\{\theta_{k},\; k=\overline{0,N-1}\bigr\}$ взаимно независимы.


Обозначим $\theta_{0,k}=\bigl\{\theta_{i},\; i=\overline{0,k}\bigr\}$,


$z_{0,k}=\bigl\{z_{i},\; i=\overline{0,k}\bigr\}$, где $z_{i}\in Z$.


Предполагаем, что условные распределения


\begin{gather*}


p_{0}(A)=P(\theta_{0}\in A),\\


%


p_{k}(A/z_{0,k-1})=


P(\theta_{k}\in A/\theta_{0,k-1}=z_{0,k-1}),\ \;


k=\overline{1,N-1},\ \; A\in \Psi,


\end{gather*}


известны.


Будем использовать обозначения:


$\xi_{0,k}=\bigl\{\xi_{i},\; i=\overline{0,k}\bigr\}$,


$L^{d}$ --- пространство квадратных матриц размера $d\times d$;


если $A\in L^{d}$, то $[A]_{ij}$ --- $(i,j)$-й элемент матрицы $A$.





Рассмотрим управляемое рекуррентное стохастическое уравнение


\begin{equation}


\xi_{k+1}=a_{k}(\theta)\xi_{k}+b_{k}(\theta)\eta_{k}


+g_{1k}(\theta,\xi_{k})w_{1k}+


g_{2k}(\theta,\eta_{k})w_{2k}+ g_{3k}(\theta)w_{3k},\ \;


k=\overline{0,N-1},


\end{equation}


где $\xi_{k}\in R^{n}$; $\eta_{k}\in R^{m}$; $\xi_{0}$ не зависит от


$(w_{1},w_{2},w_{3},\theta)$.





Предполагаем, что


$$


g_{1k}(\theta,x)=\sum_{j=1}^{n}g_{1kj}(\theta)x_{j},\quad


g_{2k}(\theta,u)=\sum_{j=1}^{m}g_{2kj}(\theta)u_{j},


$$


где $x_{j}$ и $u_{j}$ --- элементы векторов $x$ и $u$.


Коэффициенты $a_{k}(\theta)$; $b_{k}(\theta)$;


$g_{1kj}(\theta)$, $j=\overline{1,n}$;


$g_{2kj}(\theta)$, $j=\overline{1,m}$;


$g_{3k}(\theta)$, $k=\overline{0,N-1}$, --- 


$\Psi^{N}$-измеримые матричные функции соответствующих


размерностей.


Будем предполагать, что $a_{k}(\theta)=a_{k}(\theta_{0,k})$,


$b_{k}(\theta)=b_{k}(\theta_{0,k})$,


$g_{1kj}(\theta)=g_{1kj}(\theta_{0,k})$,


$g_{2kj}(\theta)=g_{2kj}(\theta_{0,k})$,


$g_{3k}(\theta)=g_{3k}(\theta_{0,k})$, т.е. коэффциценты уравнения являются


неупреждающими функционалами\linebreak от $\theta$.





Рассмотрим квадратический функционал стоимости


\begin{equation}


J(\eta)=M\biggl\{\,\sum_{k=0}^{N-1}(\xi_{k}^{*}m_{k}\xi_{k}


+\eta_{k}^{*}n_{k}\eta_{k})+\xi_{N}^{*}H\xi_{N}\biggr\}.


\end{equation}


Здесь $\bigl\{m_{k},\; k=\overline{0,N-1}\bigr\}$, 


$H$ --- симметричные неотрицательно


определенные матрицы,


$\bigl\{n_{k},\; k=\overline{0,N-1}\bigr\}$ --- симметричные


положительно определенные матрицы соответствующих размерностей.





Пусть $F_{k}=\sigma\bigl\{\xi_{0},\dotsc,\xi_{k};\theta_{0},


\dotsc,\theta_{k}\bigr\}$.


Определим класс допустимых управлений $U$ как совокупность


$\bigl\{F_{k},\; k=\overline{0,N-1}\bigr\}$-согласованных процессов


$\eta_{k}$ с $M{|\eta_{k}|}^{2}<\infty$.





Управление $\eta^{o}\in U$ будем называть оптимальным, если


$J(\eta^{o})=\inf\bigl\{J(\eta)\mid\eta\in U\bigl\}$.





Определим функции $P_{k}(\cdot):Z^{k+1}\to L^{n}$,


$R_{k}(\cdot):Z^{k+1}\to L^{n}$,


$\delta_{k}(\cdot):Z^{k+1}\to L^{n}$, $\Gamma_{k}(\cdot):Z^{k+1}\to L^{m}$


с помощью следующих рекуррентных соотношений:


\begin{gather}


P_{N}(z_{0,N})=H,\notag\\


%


R_{k}(z_{0,k})=M\bigl\{P_{k+1}(\theta_{0,k+1})/\theta_{0,k}=z_{0,k}


\bigr\}=


\int P_{k+1}(z_{0,k+1})p_{k+1}(dz_{k+1}/z_{0,k}),\notag\\


%


\bigl[\delta_{k}(z_{0,k})\bigr]_{ij}=\tr g_{1ki}^{*}(z_{0,k})R_{k}(z_{0,k})


g_{1kj}(z_{0,k}),\notag\\


%


\bigl[\Gamma_{k}(z_{0,k})\bigr]_{ij}=\tr g_{2ki}^{*}(z_{0,k})R_{k}(z_{0,k})


g_{2kj}(z_{0,k}),\notag\\


%


P_{k}(z_{0,k})=m_{k}+


a_{k}^{*}(z_{0,k})R_{k}(z_{0,k})a_{k}(z_{0,k})-


a_{k}^{*}(z_{0,k})R_{k}(z_{0,k})


b_{k}(z_{0,k})


\bigl[b_{k}^{*}(z_{0,k})R_{k}(z_{0,k})


b_{k}(z_{0,k})+\notag\\


+n_{k}+\Gamma_{k}(z_{0,k})\bigr]^{-1}


b_{k}^{*}(z_{0,k})R_{k}(z_{0,k})


a_{k}(z_{0,k})+\delta_{k}(z_{0,k}),\quad


k=\overline{0,N-1}.


\end{gather}


Отметим, что матрицы $P_{k}(z_{0,k})$,


$R_{k}(z_{0,k})$, $\delta_{k}(z_{0,k})$ и


$\Gamma_{k}(z_{0,k})$ неотрицательно определены.





\begin{lemnonum} Справедливы соотношения


\begin{align}


M\bigl\{w_{1k}^{*}g_{1k}^{*}(\theta,\xi_{k})P_{k+1}(\theta_{0,k+1})


g_{1k}(\theta,\xi_{k})w_{1k}/\xi_{0,k},\theta_{0,k}\bigr\}&= 


\xi_{k}^{*}\delta_{k}(\theta_{0,k})\xi_{k},\\


M\bigl\{w_{2k}^{*}g_{2k}^{*}(\theta,\eta_{k})P_{k+1}(\theta_{0,k+1})


g_{2k}(\theta,\eta_{k})w_{2k}/\xi_{0,k},\theta_{0,k}\bigr\}&= 


\eta_{k}^{*}\Gamma_{k}(\theta_{0,k})\eta_{k},\\


M\bigl\{w_{3k}^{*}g_{3k}^{*}(\theta)P_{k+1}(\theta_{0,k+1})


g_{3k}(\theta)w_{3k}/\xi_{0,k},\theta_{0,k}\bigr\}&= 


\tr g_{3k}^{*}(\theta)R_{k}(\theta_{0,k})g_{3k}(\theta).


\end{align}


\end{lemnonum}





\begin{pf} Имеем


\begin{gather*}


M\bigl\{w_{1k}^{*}g_{1k}^{*}(\theta,\xi_{k})P_{k+1}(\theta_{0,k+1})


g_{1k}(\theta,\xi_{k})w_{1k}/\xi_{0,k},\theta_{0,k}\bigr\}=\\


%


=M\bigl\{ \tr w_{1k}w_{1k}^{*}g_{1k}^{*}(\theta,\xi_{k})


P_{k+1}(\theta_{0,k+1})g_{1k}(\theta,\xi_{k})


/\xi_{0,k},\theta_{0,k}\bigr\}=\\


%


=M\biggl\{\tr \sum_{i,j}w_{1k}w_{1k}^{*}g_{1ki}^{*}(\theta)


P_{k+1}(\theta_{0,k+1})g_{1kj}(\theta)(\xi_{k})_{i}(\xi_{k})_{j}/


\xi_{0,k},\theta_{0,k}\biggr\}=\\


%


=\tr \sum_{i,j}M\bigl\{w_{1k}w_{1k}^{*}g_{1ki}^{*}(\theta)


P_{k+1}(\theta_{0,k+1})g_{1kj}(\theta)/


\theta_{0,k}\bigr\}(\xi_{k})_{i}(\xi_{k})_{j}.


\end{gather*}


Так как $w_{1k}$ и $\theta_{0,k}$ взаимно независимы, то


\begin{gather*}


M\bigl\{w_{1k}w_{1k}^{*}g_{1ki}^{*}(\theta)


P_{k+1}(\theta_{0,k+1})g_{1kj}(\theta)/


\theta_{0,k}\bigr\}=


M\bigl\{w_{1k}w_{1k}^{*}\bigr\}


M\bigl\{g_{1ki}^{*}(\theta)


P_{k+1}(\theta_{0,k+1})g_{1kj}(\theta)/


\theta_{0,k}\bigr\}=\\


%


=g_{1ki}^{*}(\theta_{0,k})M\bigl\{P_{k+1}(\theta_{0,k+1})/


\theta_{0,k}\bigr\}g_{1kj}(\theta_{0,k})=


g_{1ki}^{*}(\theta_{0,k})R_{k}(\theta_{0,k})g_{1kj}(\theta_{0,k}).


\end{gather*}


Таким образом,


\begin{gather*}


M\bigl\{w_{1k}^{*}g_{1k}^{*}(\theta,\xi_{k})P_{k+1}(\theta_{0,k+1})


g_{1k}(\theta,\xi_{k})w_{1k}/\xi_{0,k},\theta_{0,k}\bigr\}=\\


%


=\tr \sum_{i,j}g_{1ki}^{*}(\theta_{0,k})R_{k}


(\theta_{0,k})g_{1kj}(\theta_{0,k})


(\xi_{k})_{i}(\xi_{k})_{j}=\\


%


=\sum_{i,j}\tr g_{1ki}^{*}(\theta_{0,k})R_{k}


(\theta_{0,k})g_{1kj}(\theta_{0,k})


(\xi_{k})_{i}(\xi_{k})_{j}=


\sum_{i,j}\bigl[\delta_{k}(\theta_{0,k})\bigr]_{ij}


(\xi_{k})_{i}(\xi_{k})_{j}=


\xi_{k}^{*}\delta_{k}(\theta_{0,k})\xi_{k}.


\end{gather*}


Соотношение (4) доказано. Аналогично доказываются равенства (5) и (6).


\end{pf}





\begin{thmnonum} Оптимальное управление решением уравнения {\em (1)}


относительно функционала стоимости {\em (2)} в классе $U$ имеет вид


\begin{equation}


\eta_{k}^{o}(\theta)=-\bigl[b_{k}^{*}(\theta)R_{k}(\theta_{0,k})


b_{k}(\theta)+n_{k}+


\Gamma_{k}(\theta_{0,k})\bigr]^{-1}


b_{k}^{*}(\theta)R_{k}(\theta_{0,k})a_{k}(\theta)\xi_{k}.


\end{equation}


\end{thmnonum}





{\bf Доказательство.}


Рассмотрим выражение


\begin{gather*}


M\bigl\{\xi_{k+1}^{*}P_{k+1}(\theta_{0,k+1})\xi_{k+1}-


\xi_{k}^{*}P_{k}(\theta_{0,k})\xi_{k}\bigr\}=\\


%


=M\bigl\{M\bigl\{\xi_{k+1}^{*}P_{k+1}(\theta_{0,k+1})\xi_{k+1}-


\xi_{k}^{*}P_{k}(\theta_{0,k})\xi_{k}/


\xi_{0,k},\theta_{0,k}\bigr\}\bigr\}.


\end{gather*}


Используя представление (1) для $\xi_{k+1}$, формулу (3) и лемму,


получим


\begin{gather*}


M\bigl\{\xi_{k+1}^{*}P_{k+1}(\theta_{0,k+1})\xi_{k+1}-


\xi_{k}^{*}P_{k}(\theta_{0,k})\xi_{k}/\xi_{0,k},


\theta_{0,k}\bigr\}=\\


%


=\xi_{k}^{*}a_{k}^{*}(\theta)R_{k}(\theta_{0,k})


b_{k}(\theta)\eta_{k}+


\eta_{k}^{*}b_{k}^{*}(\theta)R_{k}(\theta_{0,k})


a_{k}(\theta)\xi_{k}+\\


%


+\eta_{k}^{*}b_{k}^{*}(\theta)R_{k}(\theta_{0,k})


b_{k}(\theta)\eta_{k}


+\eta_{k}^{*}\Gamma_{k}(\theta_{0,k})\eta_{k}+


\tr g_{3k}^{*}(\theta)R_{k}(\theta_{0,k})


g_{3k}(\theta)- \\


%


-\xi_{k}^{*}m_{k}\xi_{k}+\xi_{k}^{*}a_{k}^{*}(\theta)R_{k}(\theta_{0,k})


b_{k}(\theta)


\bigl[b_{k}^{*}(\theta)R_{k}(\theta_{0,k})


b_{k}(\theta)+n_{k}+\Gamma_{k}(\theta_{0,k})\bigr]^{-1}


b_{k}^{*}(\theta)R_{k}(\theta_{0,k})


a_{k}(\theta)\xi_{k}.


\end{gather*}


Из предыдущего равенства следует


\begin{gather*}


M\bigl\{\xi_{k+1}^{*}P_{k+1}(\theta_{0,k+1})\xi_{k+1}-


\xi_{k}^{*}P_{k}(\theta_{0,k})\xi_{k}\bigr\}=\\


%


=M\bigl\{\xi_{k}^{*}a_{k}^{*}(\theta)R_{k}(\theta_{0,k})


b_{k}(\theta)\eta_{k}+


\eta_{k}^{*}b_{k}^{*}(\theta)R_{k}(\theta_{0,k})


a_{k}(\theta)\xi_{k}+ \\


%


+\eta_{k}^{*}b_{k}^{*}(\theta)R_{k}(\theta_{0,k})


b_{k}(\theta)\eta_{k}


+\eta_{k}^{*}\Gamma_{k}(\theta_{0,k})\eta_{k}+


\tr g_{3k}^{*}(\theta)R_{k}(\theta_{0,k})


g_{3k}(\theta)- \\


%


-\xi_{k}^{*}m_{k}\xi_{k}+\xi_{k}^{*}a_{k}^{*}(\theta)R_{k}(\theta_{0,k})


b_{k}(\theta)


\bigl[b_{k}^{*}(\theta)R_{k}(\theta_{0,k})


b_{k}(\theta)+n_{k}+\Gamma_{k}(\theta_{0,k})\bigr]^{-1}


b_{k}^{*}(\theta)R_{k}(\theta_{0,k})


a_{k}(\theta)\xi_{k}\bigr\}.


\end{gather*}


Прибавим к обеим частям получившегося соотношения 


$\eta_{k}^{*}n_{k}\eta_{k}$,


перенесем влево $-\xi_{k}^{*}m_{k}\xi_{k}$ 


и просуммируем по $k$ от 0 до $N-1$.


После несложных преобразований получим


\begin{gather*}


J(\eta)=M\xi_{0}^{*}P_{0}(\theta_{0})\xi_{0}+


M\biggl\{\,\sum_{k=0}^{N-1}


\tr g_{3k}^{*}(\theta)R_{k}(\theta_{0,k})


g_{3k}(\theta)\biggr\}+ \\


%


+M\biggl\{\,\sum_{k=0}^{N-1}\bigl(\eta_{k}+\bigl[b_{k}^{*}(\theta)


R_{k}(\theta_{0,k})


b_{k}(\theta)+n_{k}+


\Gamma_{k}(\theta_{0,k})\bigr]^{-1}


b_{k}^{*}(\theta)R_{k}(\theta_{0,k})


a_{k}(\theta)\xi_{k}\bigr)^{*}


\bigl[b_{k}^{*}(\theta)R_{k}(\theta_{0,k})


b_{k}(\theta)+\\


+n_{k}+\Gamma_{k}(\theta_{0,k})\bigr]


\bigl(\eta_{k}+\bigl[b_{k}^{*}(\theta)R_{k}(\theta_{0,k})


b_{k}(\theta)+n_{k}+


\Gamma_{k}(\theta_{0,k})\bigr]^{-1}


b_{k}^{*}(\theta)R_{k}(\theta_{0,k})


a_{k}(\theta)\xi_{k}\bigr)\biggr\}.


\end{gather*}


В силу того, что матрица


$$


b_{k}^{*}(\theta)R_{k}(\theta_{0,k})


b_{k}(\theta)+n_{k}+\Gamma_{k}(\theta_{0,k})


$$


положительно определена, стоимость управления $J(\eta)$ будет минимальной


при (7), т.е.


$$


J(\eta^{o})=M\xi_{0}^{*}P_{0}(\theta_{0})\xi_{0}+


M\biggl\{\,\sum_{k=0}^{N-1}


\tr g_{3k}^{*}(\theta)R_{k}(\theta_{0,k})


g_{3k}(\theta)\biggr\}. \quad\square


$$





Рассмотрим некоторые частные случаи стохастических систем с коэффициентами,


зависящими от последовательности случайных элементов.





Пусть управляемая система описывается линейным рекуррентным стохастическим


уравнением


\begin{equation}


\xi_{k+1}=a_{k}(\theta_{k})\xi_{k}+b_{k}(\theta_{k})\eta_{k}+


g_{1k}(\theta_{k},\xi_{k})w_{1k}+


g_{2k}(\theta_{k},\eta_{k})w_{2k}+ g_{3k}(\theta_{k})w_{3k},\quad


k=\overline{0,N-1},


\end{equation}


где


$$


g_{1k}(\theta_{k},x)=\sum_{j=1}^{n}g_{1kj}(\theta_{k})x_{j},\quad


g_{2k}(\theta_{k},u)=\sum_{j=1}^{m}g_{2kj}(\theta_{k})u_{j}.


$$


Коэффициенты $a_{k}(\theta_{k})$, $b_{k}(\theta_{k})$,


$g_{1kj}(\theta_{k})$, $j=\overline{1,n}$;


$g_{2kj}(\theta_{k})$, $j=\overline{1,m}$;


$g_{3k}(\theta_{k})$, $k=\overline{0,N-1}$, --- \linebreak


$\Psi$-измеримые матричные функции соответствующих


размерностей.





\begin{cor} Оптимальное управление решением уравнения (8)


относительно функционала стоимости (2) в классе $U$ определяется


величиной $\eta^{o}_{k}(\theta_{k})$, а $\eta^{o}_{k}(\theta)$


записано в (7), причем матрицы


$R_{k}(z_{0,k})$ и $\Gamma_{k}(z_{0,k})$ находятся из рекуррентных


соотношений:


\begin{gather*}


P_{N}(z_{0,N})=H,\\


%


R_{k}(z_{0,k})=M\bigl\{P_{k+1}(\theta_{0,k+1})/\theta_{0,k}=z_{0,k}


\bigr\}=


\int P_{k+1}(z_{0,k+1})p_{k+1}(dz_{k+1}/z_{0,k}),\\


%


\bigl[\delta_{k}(z_{0,k})\bigr]_{ij}=\tr g_{1ki}^{*}(z_{k})R_{k}(z_{0,k})


g_{1kj}(z_{k}),\\


%


\bigl[\Gamma_{k}(z_{0,k})\bigr]_{ij}=\tr g_{2ki}^{*}(z_{k})R_{k}(z_{0,k})


g_{2kj}(z_{k}),\\


%


P_{k}(z_{0,k})=m_{k}+


a_{k}^{*}(z_{k})R_{k}(z_{0,k})a_{k}(z_{k})-


a_{k}^{*}(z_{k})R_{k}(z_{0,k})


b_{k}(z_{k})


\bigl[b_{k}^{*}(z_{k})R_{k}(z_{0,k})


b_{k}(z_{k})+\\


+n_{k}+\Gamma_{k}(z_{0,k})\bigr]^{-1}


b_{k}^{*}(z_{k})R_{k}(z_{0,k})


a_{k}(z_{k})+\delta_{k}(z_{0,k}),\quad k=\overline{0,N-1}.


\end{gather*}


\end{cor}





\begin{cor} Пусть $\theta_{0},\dotsc,\theta_{N-1}$ образуют 


марковскую цепь с областью


значений $D=\bigl\{1,2,\dotsc,L\bigr\}$ и переходными вероятностями


$$


p_{ij}=P(\theta_{k+1}=j/\theta_{k}=i)=


P(\theta_{k+1}=j/\theta_{0}\dotsc\theta_{k-1},\theta_{k}=i),


$$


где $i,j \in D$. Тогда оптимальное управление


решением уравнения (8) относительно функционала стоимости (2) 


в классе $U$ определяется величиной


$$


\eta_{k}^{o}=-\bigl[b_{k}^{*}(\theta_{k})R_{k}(\theta_{k})


b_{k}(\theta_{k})+n_{k}+


\Gamma_{k}(\theta_{k})\bigr]^{-1}


b_{k}^{*}(\theta_{k})R_{k}(\theta_{k})a_{k}(\theta_{k})\xi_{k},


$$


где матрицы $R_{k}(l)$ и $\Gamma_{k}(l)$, $l=\overline{1,L},$ находятся


из рекуррентных соотношений:


\begin{gather*}


P_{N}(l)=H,\\


%


R_{k}(l)=M\bigl\{P_{k+1}(\theta_{k+1})/\theta_{k}=l


\bigr\}=\sum_{j=1}^{L}P_{k+1}(j)p_{lj},\\


%


\bigl[\delta_{k}(l)\bigr]_{ij}=\tr g_{1ki}^{*}(l)R_{k}(l)


g_{1kj}(l),\\


%


\bigl[\Gamma_{k}(l)\bigr]_{ij}=\tr g_{2ki}^{*}(l)R_{k}(l)


g_{2kj}(l),\\


%


P_{k}(l)=m_{k}+


a_{k}^{*}(l)R_{k}(l)a_{k}(l)-\\


%


-a_{k}^{*}(l)R_{k}(l)


b_{k}(l)\bigl[b_{k}^{*}(l)R_{k}(l)


b_{k}(l)+n_{k}+\Gamma_{k}(l)\bigr]^{-1}


b_{k}^{*}(l)R_{k}(l)


a_{k}(l)+\delta_{k}(l),\quad k=\overline{0,N-1}.


\end{gather*}


\end{cor}





\begin{cor}


Если $\theta_{0},\dotsc,\theta_{N-1}$ --- последовательность независимых


случайных величин, то оптимальное управление


решением уравнения (8)


относительно функционала стоимости (2) в классе $U$ определяется


величиной


$$


\eta_{k}^{o}=-\bigl[b_{k}^{*}(\theta_{k})R_{k}


b_{k}(\theta_{k})+n_{k}+


\Gamma_{k}(\theta_{k})\bigr]^{-1}


b_{k}^{*}(\theta_{k})R_{k}a_{k}(\theta_{k})\xi_{k},


$$


где матрицы $R_{k}$ и $\Gamma_{k}(z)$ находятся


из рекуррентных соотношений:


\begin{gather*}


P_{N}(z)=H,\\


%


R_{k}=M\bigl\{P_{k+1}(\theta_{k+1})\bigr\}=


\int P_{k+1}(z)p_{k+1}(dz),\\


%


\bigl[\delta_{k}(z)\bigr]_{ij}=\tr g_{1ki}^{*}(z)R_{k}


g_{1kj}(z),\notag\\


%


\bigl[\Gamma_{k}(z)\bigr]_{ij}=\tr g_{2ki}^{*}(z)R_{k}


g_{2kj}(z),\\


%


P_{k}(z)=m_{k}+a_{k}^{*}(z)R_{k}a_{k}(z)-\\


%


-a_{k}^{*}(z)R_{k}


b_{k}(z)\bigl[b_{k}^{*}(z)R_{k}


b_{k}(z)+n_{k}+\Gamma_{k}(z)\bigr]^{-1}


b_{k}^{*}(z)R_{k}


a_{k}(z)+\delta_{k}(z),\quad k=\overline{0,N-1}.


\end{gather*}


\end{cor}








\begin{thebibliography}{9} 


 


\bibitem{Fra} 


Fragoso M.D. {\em Discrete-time jump LQS problem} // Int. J.


Syst. Sci. -- 1989. -- V.\,20. -- \No\,12. -- P.\,2539--2545.





\bibitem{Gi} 


Гихман И.И., Скороход А.В. {\em Управляемые случайные процессы}. --


Киев: Наукова Думка, 1977. -- 251~с.





\end{thebibliography} 


 


\vskip 2cm 


 


\post{\it Волгоградский государственный университет}{\it Поступила}


\post{}{19.04.1995} 





\end{document} 


